\preliminar{RESUMEN}
Esta memoria tiene por objetivo evaluar estrategias de cuantización de redes neuronales para seleccionar alternativas de despliegue que equilibren desempeño funcional, latencia y memoria en la plataforma STM32N6570-DK. Se consideran dos aplicaciones de inteligencia artificial en el borde: clasificación de radiografías de tórax en las categorías NORMAL y PNEUMONIA, y aproximación neuronal de un controlador predictivo basado en modelo (MPC) para un péndulo invertido sobre carro.

La metodología compara una referencia Float32 con cuantización posterior al entrenamiento INT8, entrenamiento consciente de cuantización INT8, activaciones INT16 con pesos INT8 y precisión mixta selectiva. La evaluación comprende pruebas en computador, análisis de compilación y validación en placa mediante UART, incluyendo simulación de la planta con el controlador ejecutado en el microcontrolador. Se distingue la ejecución en Cortex-M55, con soporte vectorial Helium, de la aceleración mediante la NPU Neural-ART, considerando las operaciones derivadas a CPU mediante mecanismos de fallback.

En radiografías, la cuantización posterior al entrenamiento INT8 alcanzó una exactitud de 80,93\%, frente al 79,81\% de Float32, y redujo un 59,41\% el tamaño de los archivos comparados. En MPC, las variantes INT8 disminuyeron la latencia de aproximadamente 0,303 a 0,067 ms, pero obtuvieron tasas de éxito en lazo cerrado de entre 25\% y 37\%, frente al 99,5\% de la referencia. INT16×8 conservó un 99\% de éxito en simulación; sin embargo, su despliegue requirió una conversión completa a Float32 sobre CPU, sin conservar la ejecución cuantizada nativa. La precisión mixta selectiva alcanzó un 78\% de éxito y tiempos de inferencia de entre 326 y 347 µs en las pruebas con controlador físico.

Se concluye que la estrategia de cuantización debe seleccionarse según los requisitos de cada aplicación y la representación numérica efectivamente ejecutada. PTQ INT8 permitió reducir el almacenamiento del clasificador sin deteriorar la exactitud observada. Para MPC, cuando se exige conservar operaciones cuantizadas, la precisión mixta selectiva constituye una alternativa viable entre las evaluadas, aunque su menor tasa de éxito debe contrastarse con las exigencias de control.


\textbf{Palabras clave:} inteligencia artificial en el borde; cuantización; STM32N6; Helium; Neural-ART; clasificación de radiografías; control predictivo.
